\Needspace{8\baselineskip}
\section{部署形态与完整任务的质量}

\subsection{同一任务的个人与共享参考组合}
个人工作站把模型、资料和工具安排在使用者附近，便于保留个人任务环境；集中服务则在多人之间复用模型资源，并统一更新和管理请求。表\ref{tab:reference-stack}将第\ref{sec:case}节的财报比较任务对应到两套参考设计，说明入口、模型、资料和计算工具如何组成一套系统。

\begin{table}[H]\small\centering
\caption{财报比较任务的个人与共享参考组合}\label{tab:reference-stack}
\begin{tabularx}{\linewidth}{@{}P{22mm}Y Y@{}}\toprule
组成 & 专业个人机：资料与文件助手 & 部门共享：固定财报比较应用\\\midrule
使用入口 & 本机Goose桌面入口，接收目标和文件位置。 & 内网Dify应用，提供文件、期间等固定输入项。\\
模型推理 & 专业Mac运行Gemma 4 31B IT的MLX兼容制品，由MLX-VLM提供本机推理服务。 & GPU服务器上的Qwen3.6-35B-A3B，经vLLM向应用提供推理服务。\\
解析与检索 & 本机配置PDF解析工具，保留页面位置；三份材料可直接读取。 & 应用服务器提取财报和CSV；扩大到资料库时，可另接内网嵌入与检索服务。\\
计算工具 & 本机独立任务目录中的Python计算工具，由Goose调用。 & Dify流程连接隔离计算环境，按规则完成加总、单位及比率检查。\\
结果与状态 & 比较表、说明及来源记录保存在个人任务目录，会话由助手记录。 & 结果文件保存在内部存储，应用数据库记录流程状态；复核与下载入口另行连接。\\
主要建设 & 配置本机模型、解析及计算工具，维护转换制品与任务环境。 & 配置字段、规则、导出与复核，维护身份、数据库、计算环境和共享服务。\\\bottomrule
\end{tabularx}
\source{根据Goose、MLX-VLM、Gemma、Dify及Qwen/vLLM官方能力设计，整体组合为本报告示例。\src{TECH-047}\src{TECH-024}\src{TECH-014}\src{TECH-030}\src{TECH-013}\src{TECH-020}}
\end{table}

个人组合由助手动态安排步骤，适合临时改变比较问题；共享组合提前确定处理规则，便于多位使用者重复形成同类产物。两者也可混合：Goose等个人入口保留本地文件工具，将模型请求交给集中GPU。模型运行位置和工具、资料所在位置可以分别安排。

\subsection{资源、接口与运行条件}
第\ref{sec:memory}节的内存构成可以直接对应到两种部署方式。专业个人机将模型、上下文和本机工具安排在一台设备上；集中GPU服务则通过共享权重和请求调度提高设备利用率。大型模型进入多设备或CPU/GPU异构系统后，主存容量、显存容量与设备互联共同支撑模型运行。

数十名用户可以共用模型服务，各自的数据访问、工具凭证、任务目录和结果按身份管理。这样既可以集中复用推理资源，也可以保留个人文件环境和部门资料权限。模型接口负责传递请求、输出与工具指令，应用负责执行工具、保存来源及恢复任务；具体协议对应关系见附录\ref{app:protocol}。

个人部署的持续成本包含逐台更新和设备利用差异；集中服务可以统一维护并复用资源，同时管理峰值、排队与恢复。容量分析关注实际工作时段和任务组成：注册人数决定服务覆盖范围，同时发起推理的请求数决定瞬时计算负载。

\subsection{从模型速度到业务产物}
财报比较任务把数字、引用、可复算的表格和业务解释放在同一个产物中。程序计算120到138的15\%增长，模型结合财报及注释组织说明；进一步分析增长来源时，再接入销量、价格和产品结构数据。完整任务的质量由这些环节共同形成。

\begin{table}[H]\small\centering
\caption{完整任务的评价层次}
\begin{tabularx}{\linewidth}{@{}P{29mm}Y Y@{}}\toprule
层次 & 核心问题 & 可记录的结果\\\midrule
模型内容 & 事实、数字、事件状态和引用是否正确。 & 准确性、遗漏、冲突、恰当弃答与修订量。\\
推理服务 & 是否在可接受等待内完成请求。 & 首次响应、完整时延、达标吞吐与失败率。\\
应用任务 & 工具和资料能否形成所需产物。 & 任务完成、文件可用性、恢复与人工介入。\\
持续运行 & 用户与运维能否稳定使用。 & 状态恢复、版本变化、资源和维护时间。\\\bottomrule
\end{tabularx}
\end{table}

设备支出决定可用资源，模型调用与维护影响持续成本，人工修订则决定节省时间能否兑现。由此，经济性可以围绕完成一项合格任务所用的机器资源、人员时间和维护投入展开，将推理性能与实际工作效率联系起来。

\Needspace{10\baselineskip}
\section{阶段认识}
当前开放模型已为专业个人工作站和集中服务提供多种能力基础。Qwen与Gemma覆盖长文、多模态和工具使用，稠密与MoE结构提供不同的计算和存储配置；DeepSeek、GLM和Kimi进一步面向复杂研究、办公文件与长过程任务。模型能力、每次计算量和整体内存需求是理解这些路线的三个主要维度。

推理服务把模型能力转化为持续可调用的计算资源。vLLM和SGLang通过调度、缓存与并行服务多人，MLX利用Apple设备的统一内存，异构路线利用CPU和主存扩展模型装载。新模型的专用适配与正式发布包形成不同的更新节奏，对应版本固定、升级和维护的具体工作。

非编码应用则决定同事如何完成任务。财报比较例子中，工作台支持边读边问，Dify将方法固化为重复流程，Goose围绕目标安排文件和计算操作。开发框架进一步连接专用数据与内部流程。这些应用形态已经把模型使用从问答扩展到资料搜集、计算、文件生成和连续任务。

个人环境提供任务连续性和本机操作的便利，共享服务提供资源复用与集中维护，两者可以组合使用。经济性的核心是完成一项工作所消耗的设备资源、人员时间与维护投入：模型提高内容处理能力，服务提高计算利用率，应用减少资料和工具之间的人工衔接，三层共同形成业务效率。
